\chapter{Introducción al Lenguaje C}

\section{Del diagrama de flujo al código}

En el bloque anterior aprendimos a representar algoritmos con diagramas de flujo. Ahora vamos a traducir esos diagramas a un lenguaje que la computadora puede compilar y ejecutar: el \textbf{lenguaje C}.

C es un lenguaje de propósito general creado en los años 70. Es el lenguaje más usado en sistemas embebidos y microcontroladores por dos razones principales: es \textbf{eficiente} (genera código compacto y rápido) y ofrece \textbf{acceso directo al hardware} (registros, puertos, memoria). Todo lo que aprendamos en este bloque se va a aplicar directamente en la programación del microcontrolador.

\section{Estructura de un programa en C}

Todo programa en C tiene la siguiente estructura mínima:

\begin{lstlisting}[style=estilo-c]
#include <stdio.h>      // Incluir biblioteca estándar de entrada/salida

int main(void) {        // Función principal — el programa empieza aquí
    // instrucciones
    return 0;           // Indicar que el programa terminó correctamente
}
\end{lstlisting}

\textbf{Elementos obligatorios:}

\begin{itemize}
    \item \texttt{\#include <stdio.h>}: incluye la biblioteca que permite usar \texttt{printf} y \texttt{scanf} (entrada y salida por consola). En programas para microcontrolador esto cambiará, pero por ahora la usamos para probar en la PC.
    \item \texttt{int main(void)}: función principal. Todo programa en C debe tener exactamente una función \texttt{main}. La ejecución empieza en la primera línea de \texttt{main} y termina cuando se llega al \texttt{return}.
    \item \texttt{\{ \}}: las llaves definen un \textbf{bloque} de código. Todo lo que está entre \texttt{\{} y \texttt{\}} pertenece a ese bloque.
    \item \texttt{;}: cada instrucción termina con punto y coma. Olvidarlo es uno de los errores más frecuentes.
    \item \texttt{// comentario}: el compilador ignora todo lo que está a la derecha de \texttt{//} en esa línea. Los comentarios sirven para explicar el código.
\end{itemize}

\section{Variables y tipos de datos}

Antes de usar una variable en C hay que \textbf{declararla}: indicar su nombre y el tipo de dato que va a almacenar.

\begin{lstlisting}[style=estilo-c]
tipo nombre;
tipo nombre = valor_inicial;
\end{lstlisting}

\subsection{Tipos de datos básicos}

\begin{table}[H]
\centering
\begin{tabularx}{\textwidth}{ll>{\raggedright\arraybackslash}X>{\raggedright\arraybackslash}X}
\toprule
Tipo & Tamaño típico & Rango aproximado & Uso \\
\midrule
\texttt{int} & 2 o 4 bytes & $-32.768$ a $32.767$ (2 bytes) & Números enteros \\
\texttt{unsigned int} & 2 o 4 bytes & 0 a 65.535 (2 bytes) & Enteros sin signo \\
\texttt{float} & 4 bytes & $\pm 3{,}4 \times 10^{38}$ & Números decimales (simple precisión) \\
\texttt{double} & 8 bytes & $\pm 1{,}7 \times 10^{308}$ & Números decimales (doble precisión) \\
\texttt{char} & 1 byte & $-128$ a $127$ & Un carácter ASCII o un entero de 8 bits \\
\texttt{unsigned char} & 1 byte & 0 a 255 & Byte sin signo --- muy usado en embebidos \\
\bottomrule
\end{tabularx}
\end{table}

\begin{notabox}
En programación de microcontroladores AVR, el tipo más usado es \texttt{uint8\_t} (equivalente a \texttt{unsigned char}, siempre 8 bits). Por ahora trabajamos con \texttt{int} y \texttt{float} para mayor comodidad.
\end{notabox}

\subsection{Ejemplos de declaración}

\begin{lstlisting}[style=estilo-c]
int edad;               // declara una variable entera sin valor inicial
int contador = 0;       // declara e inicializa en 0
float temperatura = 25.5;
char letra = 'A';
unsigned int suma = 0;
\end{lstlisting}

\subsection{Reglas para nombrar variables}

\begin{itemize}
    \item Solo letras, números y guion bajo \texttt{\_}. No puede empezar con un número.
    \item C distingue mayúsculas de minúsculas: \texttt{Temperatura} y \texttt{temperatura} son variables distintas.
    \item No se pueden usar palabras reservadas del lenguaje (\texttt{int}, \texttt{if}, \texttt{while}, \texttt{return}, etc.).
    \item Elegir nombres descriptivos: \texttt{temperatura} es mejor que \texttt{t} o \texttt{x}.
\end{itemize}

\section{Operadores}

\subsection{Aritméticos}

\begin{table}[H]
\centering
\begin{tabular}{cccc}
\toprule
Operador & Operación & Ejemplo & Resultado \\
\midrule
\texttt{+} & Suma & \texttt{3 + 2} & \texttt{5} \\
\texttt{-} & Resta & \texttt{5 - 3} & \texttt{2} \\
\texttt{*} & Multiplicación & \texttt{4 * 3} & \texttt{12} \\
\texttt{/} & División & \texttt{7 / 2} & \texttt{3} (entera si ambos son \texttt{int}) \\
\texttt{\%} & Módulo (resto) & \texttt{7 \% 2} & \texttt{1} \\
\bottomrule
\end{tabular}
\end{table}

\begin{advertenciabox}
\textbf{Atención con la división entera:} en C, dividir dos \texttt{int} descarta los decimales. \texttt{7 / 2} da \texttt{3}, no \texttt{3.5}. Para obtener el decimal hay que usar \texttt{float}: \texttt{7.0 / 2} da \texttt{3.5}.
\end{advertenciabox}

\subsection{Relacionales}

\begin{table}[H]
\centering
\begin{tabular}{cl}
\toprule
Operador & Significado \\
\midrule
\texttt{==} & Igual (doble signo --- no confundir con \texttt{=} de asignación) \\
\texttt{!=} & Distinto \\
\texttt{>} & Mayor que \\
\texttt{<} & Menor que \\
\texttt{>=} & Mayor o igual \\
\texttt{<=} & Menor o igual \\
\bottomrule
\end{tabular}
\end{table}

\subsection{Lógicos}

\begin{table}[H]
\centering
\begin{tabular}{cl}
\toprule
Operador & Significado \\
\midrule
\texttt{\&\&} & AND lógico \\
\texttt{||} & OR lógico \\
\texttt{!} & NOT lógico \\
\bottomrule
\end{tabular}
\end{table}

\subsection{Asignación}

\begin{lstlisting}[style=estilo-c]
x = 5;          // asigna 5 a x
x = x + 1;      // incrementa x en 1
x += 1;         // equivalente al anterior (forma abreviada)
x++;            // también equivalente — incremento en 1
x--;            // decremento en 1
\end{lstlisting}

\section{Entrada y salida por consola}

Para los ejercicios en PC usamos dos funciones de \texttt{<stdio.h>}:

\subsection{\texttt{printf} --- mostrar datos en pantalla}

\begin{lstlisting}[style=estilo-c]
printf("Hola mundo\n");             // imprime texto
printf("El resultado es %d\n", x); // imprime texto + valor de x (entero)
printf("Temperatura: %.2f\n", t);  // imprime t con 2 decimales
\end{lstlisting}

Los especificadores de formato más comunes:

\begin{table}[H]
\centering
\begin{tabular}{cl}
\toprule
Especificador & Tipo \\
\midrule
\texttt{\%d} & Entero (\texttt{int}) \\
\texttt{\%f} & Decimal (\texttt{float} o \texttt{double}) \\
\texttt{\%.2f} & Decimal con 2 cifras decimales \\
\texttt{\%c} & Carácter (\texttt{char}) \\
\texttt{\%s} & Cadena de texto \\
\bottomrule
\end{tabular}
\end{table}

El \texttt{\textbackslash n} al final provoca un salto de línea.

\subsection{\texttt{scanf} --- leer datos del teclado}

\begin{lstlisting}[style=estilo-c]
int x;
scanf("%d", &x);    // lee un entero y lo guarda en x

float t;
scanf("%f", &t);    // lee un decimal y lo guarda en t
\end{lstlisting}

\begin{notabox}
El \texttt{\&} antes del nombre de la variable es obligatorio en \texttt{scanf}. Indica la \textbf{dirección de memoria} donde se debe guardar el dato. Olvidarlo causa comportamiento impredecible.
\end{notabox}

\section{Estructuras de control}

\subsection{Secuencia}

Las instrucciones se ejecutan en orden, de arriba a abajo. No requiere ninguna sintaxis especial.

\begin{lstlisting}[style=estilo-c]
int a, b, resultado;
scanf("%d", &a);
scanf("%d", &b);
resultado = a + b;
printf("Resultado: %d\n", resultado);
\end{lstlisting}

\subsection{Decisión: \texttt{if} / \texttt{else}}

\begin{lstlisting}[style=estilo-c]
if (condición) {
    // se ejecuta si la condición es verdadera
} else {
    // se ejecuta si la condición es falsa (opcional)
}
\end{lstlisting}

Para condiciones encadenadas:

\begin{lstlisting}[style=estilo-c]
if (temperatura > 26) {
    printf("Hace calor\n");
} else if (temperatura < 22) {
    printf("Hace frio\n");
} else {
    printf("Temperatura agradable\n");
}
\end{lstlisting}

\begin{advertenciabox}
\textbf{Errores frecuentes:}
\begin{itemize}
    \item Usar \texttt{=} en lugar de \texttt{==} dentro del \texttt{if}: \texttt{if (x = 5)} asigna 5 a x y siempre es verdadero. Lo correcto es \texttt{if (x == 5)}.
    \item Olvidar las llaves \texttt{\{\}} cuando el bloque tiene más de una instrucción.
\end{itemize}
\end{advertenciabox}

\subsection{Repetición: \texttt{while}}

La condición se evalúa \textbf{antes} de cada iteración. Si es falsa desde el principio, el cuerpo no se ejecuta.

\begin{lstlisting}[style=estilo-c]
while (condición) {
    // cuerpo del bucle
}
\end{lstlisting}

\begin{lstlisting}[style=estilo-c]
int contador = 1;
while (contador <= 10) {
    printf("%d\n", contador);
    contador++;
}
\end{lstlisting}

\subsection{Repetición: \texttt{do-while}}

El cuerpo se ejecuta \textbf{al menos una vez}; la condición se evalúa al final.

\begin{lstlisting}[style=estilo-c]
do {
    // cuerpo del bucle
} while (condición);
\end{lstlisting}

\begin{lstlisting}[style=estilo-c]
int valor;
do {
    printf("Ingrese un numero positivo: ");
    scanf("%d", &valor);
} while (valor <= 0);
\end{lstlisting}

\subsection{Repetición: \texttt{for}}

Forma compacta de escribir un bucle cuando se conoce de antemano la cantidad de iteraciones.

\begin{lstlisting}[style=estilo-c]
for (inicialización; condición; actualización) {
    // cuerpo del bucle
}
\end{lstlisting}

\begin{lstlisting}[style=estilo-c]
for (int i = 1; i <= 10; i++) {
    printf("%d\n", i);
}
\end{lstlisting}

Equivale exactamente a:

\begin{lstlisting}[style=estilo-c]
int i = 1;
while (i <= 10) {
    printf("%d\n", i);
    i++;
}
\end{lstlisting}

\section{Equivalencia entre pseudocódigo y C}

\begin{table}[H]
\centering
\begin{tabular}{ll}
\toprule
Pseudocódigo (PSeInt) & C \\
\midrule
\texttt{Leer variable} & \texttt{scanf("\%d", \&variable);} \\
\texttt{Escribir "texto"} & \texttt{printf("texto\textbackslash n");} \\
\texttt{Escribir variable} & \texttt{printf("\%d\textbackslash n", variable);} \\
\texttt{variable \textleftarrow\ valor} & \texttt{variable = valor;} \\
\texttt{Si condición Entonces ... Fin Si} & \texttt{if (condición) \{ ... \}} \\
\texttt{Si ... Sino ... Fin Si} & \texttt{if (...) \{ ... \} else \{ ... \}} \\
\texttt{Mientras condición Hacer ... Fin Mientras} & \texttt{while (condición) \{ ... \}} \\
\texttt{Repetir ... Hasta Que condición} & \texttt{do \{ ... \} while (condición);} \\
\texttt{MOD} & \texttt{\%} \\
\texttt{Y} & \texttt{\&\&} \\
\texttt{O} & \texttt{||} \\
\texttt{NO} & \texttt{!} \\
\bottomrule
\end{tabular}
\end{table}

\section{Ejercicios: los algoritmos del TP en C}

Los siguientes programas son la traducción directa de los diagramas de flujo del Bloque 4. Se pueden compilar y ejecutar en cualquier compilador de C estándar (online: repl.it, OnlineGDB, Tutorialspoint; local: VSCode con GCC).

\subsection{Ejercicio 1 --- Suma de dos valores}

\begin{lstlisting}[style=estilo-c]
#include <stdio.h>

int main(void) {
    int a, b, resultado;

    printf("Ingrese el primer numero: ");
    scanf("%d", &a);

    printf("Ingrese el segundo numero: ");
    scanf("%d", &b);

    resultado = a + b;

    printf("La suma es: %d\n", resultado);

    return 0;
}
\end{lstlisting}

\subsection{Ejercicio 2 --- Operación con tres valores}

\begin{lstlisting}[style=estilo-c]
#include <stdio.h>

int main(void) {
    int a, b, c, resultado;

    printf("Ingrese el primer numero: ");
    scanf("%d", &a);

    printf("Ingrese el segundo numero: ");
    scanf("%d", &b);

    printf("Ingrese el tercer numero: ");
    scanf("%d", &c);

    resultado = (a + b) - c;

    printf("El resultado es: %d\n", resultado);

    return 0;
}
\end{lstlisting}

\subsection{Ejercicio 3 --- Módulo (valor absoluto)}

\begin{lstlisting}[style=estilo-c]
#include <stdio.h>

int main(void) {
    int numero, modulo;

    printf("Ingrese un numero: ");
    scanf("%d", &numero);

    if (numero < 0) {
        modulo = -numero;
    } else {
        modulo = numero;
    }

    printf("El modulo es: %d\n", modulo);

    return 0;
}
\end{lstlisting}

\begin{notabox}
La biblioteca estándar \texttt{<math.h>} tiene la función \texttt{abs()} que hace esto automáticamente, pero aquí lo implementamos manualmente para practicar la estructura \texttt{if/else}.
\end{notabox}

\subsection{Ejercicio 4 --- Sensor de temperatura}

\begin{lstlisting}[style=estilo-c]
#include <stdio.h>

int main(void) {
    float temperatura;

    printf("Ingrese la temperatura: ");
    scanf("%f", &temperatura);

    if (temperatura > 26.0) {
        printf("Hace calor\n");
    } else if (temperatura < 22.0) {
        printf("Hace frio\n");
    }

    return 0;
}
\end{lstlisting}

\begin{notabox}
Cuando la temperatura está entre 22° y 26° el programa no imprime nada, tal como pide el enunciado.
\end{notabox}

\subsection{Ejercicio 5 --- Promedio de una serie de valores}

\begin{lstlisting}[style=estilo-c]
#include <stdio.h>

int main(void) {
    int numero;
    float suma = 0;
    int cantidad = 0;

    printf("Ingrese numeros (0 para terminar):\n");

    scanf("%d", &numero);

    while (numero != 0) {
        suma = suma + numero;
        cantidad++;
        scanf("%d", &numero);
    }

    if (cantidad > 0) {
        printf("El promedio es: %.2f\n", suma / cantidad);
    } else {
        printf("No se ingresaron numeros.\n");
    }

    return 0;
}
\end{lstlisting}

\begin{notabox}
\textbf{Puntos importantes de este ejercicio:}
\begin{itemize}
    \item La suma se declara como \texttt{float} para que la división no pierda los decimales.
    \item El \texttt{if (cantidad > 0)} evita la división por cero si el usuario ingresa 0 inmediatamente.
    \item El 0 final \textbf{no se suma} porque la condición del \texttt{while} se evalúa antes de acumular.
\end{itemize}
\end{notabox}

\section{Entorno de desarrollo}

Para compilar y ejecutar programas en C se puede usar cualquiera de los siguientes entornos:

\textbf{Online (sin instalación):}
\begin{itemize}
    \item \texttt{replit.com} --- crea un proyecto de tipo C.
    \item \texttt{onlinegdb.com} --- compilador + depurador online.
    \item \texttt{tutorialspoint.com/compile\_c\_online.php}
\end{itemize}

\textbf{Local:}
\begin{itemize}
    \item \textbf{VSCode} + extensión C/C++ + compilador GCC (a través de MinGW en Windows).
\end{itemize}
